\clearpage
\section{热力学补充}
\label{chap:appendix-thermodynamics-supplement}
\label{app:thermodynamics-supplement}
\renewcommand{\thestatement}{B.1.\arabic{statement}}
\renewcommand{\theHstatement}{appB.sec01.\arabic{statement}}
\setcounter{statement}{0}
\renewcommand{\theequation}{B.\arabic{equation}}
\renewcommand{\theHequation}{appB.\arabic{equation}}
\setcounter{equation}{0}

本附录的目标是以极其简洁, 因而必然非常不完整的方式, 回顾 Chapter~\ref{chap:fluid-mechanics-equations} 中使用的流体热力学主要概念. 我们的基本目标是引入熵的概念和热力学第二定律 (特别是 Clausius 定理). 在关于这一主题的各种参考教材中, 例如可参阅文献\MBUnresolvedReference{bib:thermodynamics-reference}{[94]}以获得远为详细的论述.

热力学研究物质系统及其能量的演化, 并考虑其内部的分子结构. 这种物质系统可以用若干称为状态变量的量很好地描述, 其中包括:
\begin{itemize}
 \item 体积, 质量和密度, 它们描述所研究系统的大小和系统中的物质量;
 \item 温度和能量, 它们描述物质基本粒子的激发状态;
 \item 压力, 它通过粒子在置于物质中的一个 (虚拟) 壁面上的碰撞, 对粒子的随机运动给出一种力学度量.
\end{itemize}
处理复杂系统 (例如多相系统) 时, 往往还需要其他状态变量.

在所有这些量中, 有些量在系统大小增加时并不改变 (温度, 压力等), 它们称为强度变量. 相反, 有些量随系统大小增加而增加 (体积, 能量等), 它们称为广延变量.

事实上, 状态变量的值并不相互独立. 通常认为, 每个均匀热力学系统都满足一个状态方程, 它是所有状态变量之间的代数关系. 根据所研究的系统, 这个方程可以写成压力关于系统温度和密度的表达式, 也可以写成温度关于压力和密度的表达式, 或者写成独立变量的任何其他组合.

在流体动力学情形中, 需要把本附录的概念应用于流体微团 $(\Omega_t)_t$. 假定它们处于瞬时局部热力学平衡, 从而每个基本流体粒子都可以视为一个封闭的均匀热力学系统. 这意味着, 在流体中任意点 $x$ 和任意时刻 $t$, 趋向局部分子平衡的弛豫时间都远短于流体整体演化的任何特征时间.

\subsection{热容量}
\label{sec:app-b-heat-capacity}

\begin{Definition}
\label{def:app-b-heat-capacity}
定容热容量 (相应地, 定压热容量) 记为 $C_V$ (相应地记为 $C_P$), 它是指在保持体积 (相应地, 压力) 恒定时, 使系统温度升高 $1\,\mathrm K$ 所需的热量. 因而, 要使温度增加 $\Delta T$, 必须向系统提供如下能量:
\[
 \Delta Q=C_V\Delta T,\qquad\text{当 }V\text{ 恒定时},
\]
\[
 \Delta Q=C_P\Delta T,\qquad\text{当 }P\text{ 恒定时}.
\]
当然, $C_V$ 和 $C_P$ 都是广延量, 因此常使用如下定义的比热容 (即单位质量的热容量):
\[
 c_V=\frac{C_V}{M},\qquad c_P=\frac{C_P}{M},
\]
其中 $M$ 是所考虑系统的质量.
\end{Definition}
